The key finding of this study was that no single generative model dominated all quality dimensions: GARCH(1,1) better reproduced volatility clustering but failed distributional tests, the standard HMM passed distributional tests but could not generate persistent high-volatility regimes, and the hybrid framework, which augmented quantile-defined regime switching with a Poisson jump-duration mechanism, delivered the most balanced performance by achieving high distributional pass rates both in-sample and out-of-sample with partial reproduction of volatility clustering governed by two tunable hyperparameters. Direct frequentist estimation eliminated the computational cost and initialization sensitivity of the EM algorithm, scaling the approach to a 424-asset pipeline, and replacing the Single-Index Model's linear factor structure with a Student-t copula substantially improved per-asset distributional accuracy for assets that did not closely track the index. Several directions for future work follow from these results: time-varying transition matrices, estimated via rolling windows or Bayesian online learning, would allow the model to adapt to structural regime shifts; scaling the Student-t copula to larger asset universes via factor copulas or truncated vine structures would extend the per-asset accuracy gains to full portfolio applications; and embedding the hybrid HMM within a portfolio optimization loop, where synthetic scenarios drive tail-risk objectives, would provide an end-to-end test of the framework's practical value.
